% Quantitative choice of atom budgets and uniform scale thresholds.
\section{Choosing the local paid budgets}
\label{sec:ex-lattice-local-paid-budgets}

The one-step insertion uses two kinds of local input: a finite list of
coefficient-recovery probes, and positive budgets for the basis, regular,
and forced atoms. The first choice is elementary once the probe list is
finite. The eventual threshold can then be made monotone in the scale and
uniform over all admissible adjacent levels. This section proves those two
selection steps. The local atom estimates themselves are recorded in the
source audit; they are the remaining inputs to the one-step recurrence of
Section~\ref{sec:ex-lattice-trace-price-recurrence}.

Fix a finite set \(\mathcal T\) of probes and nonnegative constants
\(O,H,W,G_d,a,c\), where \(O\) is the motion-average norm factor, \(H\) the local
readout factor, \(W\) the test factor, \(G_d\) the gauge-variation factor,
\(a\) the block weight, and \(c\) the tail-decay factor. For each probe
\(\tau\), let \(\pi_\tau\geq0\) be its variation price and define its
local coefficient
\[
 q_\tau=4aG_d\pi_\tau+2a.
\]
Given a target price \(y>0\), the required atom budget inequalities are
\begin{align}
 q_\tau v&\leq u &&(\tau\in\mathcal T), \label{eq:ex-lattice-atom-cycle}\\
 O\bigl(Wu+2Hav\bigr)&\leq y, \label{eq:ex-lattice-atom-regular}\\
 4OHav&\leq y, \label{eq:ex-lattice-atom-large}\\
 O\bigl(Wv+2Havc\bigr)&\leq y. \label{eq:ex-lattice-atom-main}
\end{align}
The two outputs \(u,v\) are the probe and summability budgets. In the
Wilson application, \(q_\tau\) bounds the local variation coefficient,
while the other three inequalities pay respectively for regular terms,
large-support terms, and the primary atom.

\begin{lemma}[Positive local budget choice]
\label{lem:ex-lattice-local-budget-choice}
For every finite probe family and every \(y>0\), there are \(u,v>0\)
satisfying \eqref{eq:ex-lattice-atom-cycle}--\eqref{eq:ex-lattice-atom-main}.
\end{lemma}
\begin{proof}
Set
\[
 K=\sum_{\tau\in\mathcal T}q_\tau,\qquad
 B=\max\{1,\ O[W(1+K)+2Ha],\ 4OHa,\ O[W+2Hac]\}.
\]
These are finite nonnegative numbers. Choose
\[
 v=\frac{y}{2B},\qquad u=(1+K)v.
\]
Both are strictly positive. Since \(q_\tau\leq K\leq1+K\),
\(q_\tau v\leq u\). The left side of the regular-term inequality is
\(vO[W(1+K)+2Ha]\leq vB=y/2\). The large-support and primary left
sides are respectively \(4OHav\leq Bv=y/2\) and
\(vO[W+2Hac]\leq Bv=y/2\). Each is therefore at most \(y\), as
required.
\end{proof}

We next record the uniform threshold selection used to turn fixed-parameter
estimates into a single scale profile. At a scale index \(j\), the allowed
adjacent indices are the same-level index and its immediate predecessor;
there are at most two. The coefficient grades to be checked form the finite
set \(\{0,\ldots,j-1\}\). Thus only finitely many eventual requirements
occur at each fixed \(j\).

\begin{lemma}[Monotone threshold envelope]
\label{lem:ex-lattice-monotone-threshold}
Let \(b_j\in\mathbb R\) be any prescribed lower threshold. For each
\(j\), let \(S_j\) be a finite set and, for each \(s\in S_j\), let
\(E_{j,s}\in\mathbb R\) have the property that a desired conclusion
\(\mathcal P_{j,s}(x)\) holds whenever \(x\geq E_{j,s}\). Then there is a
nondecreasing sequence \(B_j\) such that \(B_j\geq b_j\) and, for every
\(s\in S_j\), \(\mathcal P_{j,s}(x)\) holds whenever \(x\geq B_j\).
\end{lemma}
\begin{proof}
Define
\[
 e_j=\max\bigl(\{b_j\}\cup\{E_{j,s}:s\in S_j\}\bigr),
 \qquad B_j=\max_{0\leq r\leq j}e_r.
\]
The maxima exist because the sets are finite and nonempty after adjoining
\(b_j\). The prefix maximum \(B_j\) is nondecreasing and dominates
\(b_j\). If \(s\in S_j\), then \(E_{j,s}\leq e_j\leq B_j\), so every
\(x\geq B_j\) satisfies the stipulated conclusion.
\end{proof}

Apply this lemma at each scale to the finitely many pairs consisting of an
admissible adjacent level and a coefficient grade below that scale. The
atom-budget lemma supplies positive budgets for each fixed pair; the
finite-volume local estimates then supply an eventual threshold for the
corresponding primary, regular, and forced atom conclusions. Taking the
finite maximum and then the prefix maximum gives one monotone threshold
profile that works simultaneously for every adjacent transition and every
permitted coefficient grade. This is the uniformization step needed before
the scale induction. It does not prove the local atom estimates: their
mathematical inputs and exact source route are itemized in the accompanying
evidence file.
